\documentclass{article}
\usepackage{amsfonts}
\usepackage{amssymb}
\usepackage{amsmath}
\usepackage{mathtools}
\usepackage{mathrsfs}
%\usepackage{bbm}
\usepackage{xcolor}
\usepackage[utf8]{inputenc}
\usepackage[margin=1.5in]{geometry}
\usepackage[T1]{fontenc}
\usepackage{amsthm}
\usepackage[english]{babel}
\usepackage{graphicx}
\usepackage{tikz-cd}
\usepackage{tikz}
\usepackage{rotating}
%\usepackage{dsfont}
\usepackage{esint}
\usetikzlibrary{matrix, arrows}

\newtheorem{theorem}{Theorem}[subsection]
\newtheorem{prop}[theorem]{Proposition}
\newtheorem{conjecture}[theorem]{Conjecture}
\newtheorem{heuristic}[theorem]{Heuristic}
\newtheorem{lemma}[theorem]{Lemma}
\newtheorem{corollary}[theorem]{Corollary}
\theoremstyle{definition}
\newtheorem{example}[theorem]{Example}
\theoremstyle{definition}
\newtheorem{remark}[theorem]{Remark}
\theoremstyle{definition}
\newtheorem{defi}[theorem]{Definition}
\newtheorem*{theoremstar}{Theorem}

% typesetting macros
\newcommand\rea[1]{\textcolor{orange}{#1}}
\newcommand\bsni{\bigskip\noindent}

% general math macros
\newcommand\bba{\mathbb{A}}
\newcommand\bbb{\mathbb{B}}
\newcommand\bbc{\mathbb{C}}
\newcommand\bbd{\mathbb{D}}
\newcommand\bbf{\mathbb{F}}
\newcommand\bbl{\mathbb{L}}
\newcommand\bbn{\mathbb{N}}
\newcommand\bbp{\mathbb{P}}
\newcommand\bbq{\mathbb{Q}}
\newcommand\bbr{\mathbb{R}}
\newcommand\bbs{\mathbb{S}}
\newcommand\cA{\mathcal{A}}
\newcommand\cB{\mathcal{B}}
\newcommand\cC{\mathcal{C}}
\newcommand\cD{\mathcal{D}}
\newcommand\cE{\mathcal{E}}
\newcommand\cF{\mathcal{F}}
\newcommand\cG{\mathcal{G}}
\newcommand\cH{\mathcal{H}}
\newcommand\cI{\mathcal{I}}
\newcommand\cJ{\mathcal{J}}
\newcommand\cL{\mathcal{L}}
\newcommand\cM{\mathcal{M}}
\newcommand\cN{\mathcal{N}}
\newcommand\cO{\mathcal{O}}
\newcommand\cP{\mathcal{P}}
\newcommand\cR{\mathcal{R}}
\newcommand\cT{\mathcal{T}}
\newcommand\cX{\mathcal{X}}
\newcommand\cY{\mathcal{Y}}
\newcommand\cZ{\mathcal{Z}}
\newcommand\K{\mathrm{K}}
\newcommand\R{\mathrm{R}}
\newcommand\mi{^{-1}}
\newcommand\minfty{{-\infty}}
\newcommand\an{^{\mathrm{an}}}
\newcommand\red{\mathrm{red}}
\newcommand\tmax{_\mathrm{tmax}}
\newcommand\sk{\mathrm{Sk}}
\newcommand\pdisc{\overline\bbd^*}
\newcommand\disc{\overline\bbd}
\newcommand\Sing{\mathrm{Sing}}

% general AG macros
\DeclareMathOperator{\ord}{ord}
\DeclareMathOperator{\exc}{exc}
\DeclareMathOperator{\vol}{vol}
\DeclareMathOperator{\Specop}{Spec}
\DeclareMathOperator{\Projop}{Proj}
\newcommand\Spec[1]{\Specop \mathrm{#1}}
\newcommand\Proj[1]{\Projop \mathrm{#1}}
\newcommand\charac[1]{\mathrm{char}\, #1}

% n-A PSH macros
\newcommand\field{\mathrm{K}}
\newcommand\fieldext{\mathfrak{L}}
\newcommand\valring{{\field^\circ}}
\newcommand\resfield{{\tilde\field}}
\newcommand\shilov{^{\mathrm{Shi}}}
\newcommand\norm{\zeta}
\newcommand\metric{\phi}
\newcommand\PSH{\mathrm{PSH}}
\newcommand\MA{\,\mathrm{MA}}
\newcommand\grnorm{\norm_\bullet}
\newcommand\grnormspace{\mathcal{N}_\bullet}
\newcommand\normspace{\mathcal{N}}
\newcommand\gm{\mathbb{G}_m}
\newcommand\basis{\mathfrak{s}}
\newcommand\basislong{\mathfrak{s}=(s_i)_i}
\newcommand\refmetric{\phi_{\mathrm{ref}}}
\newcommand\na{^{\mathrm{NA}}}


% paper-specific macros or other
\newcommand\anaff{\mathbf{A}}
\newcommand\vp{\varphi}
\newcommand\nv{\mathcal{N}(V)}
\newcommand\cpe{\mathrm{(ce)}}
\newcommand\nlin{\mathcal{N}^{lin}_\bullet(L)}
\newcommand\ep{\mathcal{E}^p}
\newcommand\pathspace{\mathcal{P}}
\newcommand\path{\mathfrak{P}}
\newcommand\rayspace{\mathcal{R}}
\newcommand\ray{\mathfrak{R}}
\newcommand\inj{\imath}
\newcommand\raynorm{\mathfrak{r}^c}
\newcommand\raymetric{\mathfrak{R}^c}
\newcommand\pright{\dot{\phi}_0{}}
\newcommand\pleft{\dot{\phi}_1{}}
\newcommand\laurent{{\bbc((z))}}
\newcommand\power{{\bbc[[z]]}}
\newcommand\ideal{\mathfrak{a}}
\newcommand\xlaurent{X_\laurent}
\newcommand\xanlaurent{X^{\mathrm{an}}_\laurent}
\newcommand\cxpower{\cX_\power}
\newcommand\annu{{\mathrm{a}(0;a,b)}}
\newcommand\cannu{{\bar{\mathrm{a}}(0;a,b)}}
\newcommand\FS{\mathrm{FS}}
\newcommand\Gm{\mathbb{G}_\mathrm{m}}

\author{REBOULET Rémi}
\title{CURRICULUM VITAE.}
\begin{document}

%\begin{center}
%\large{\textbf{CURRICULUM VITAE.}}
%\end{center}

%\par\noindent\rule{\textwidth}{0.4pt}
%\textbf{INFORMATIONS PERSONNELLES.}

%\bsni \textbf{Nom :} REBOULET\\
%\textbf{Prénom :} Rémi\\
%\textbf{Date de naissance :} 26 Février 1997\\
%\textbf{Adresse e-mail (pérenne) :} \verb|rebouletremimath@gmail.com|\\
%\textbf{Adresse e-mail (actuelle) :} \verb|reboulet@math.univ-lyon1.fr|

%\bsni

\iffalse \par\noindent\rule{\textwidth}{0.4pt}
\textbf{INFORMATIONS PROFESSIONNELLES.}

\bsni \textbf{Octobre 2024.} Chargé de recherches (CNRS, section 41).\\
Rattaché à l'Institut Camille Jordan, UMR5208.

\bsni \textbf{Octobre 2022 - Septembre 2024 :} Postdoc en mathématiques.\\
Collaboration avec Robert BERMAN.\\
Mathematical Sciences, Chalmers University of Technology, Göteborg, Suède.

\bsni \textbf{Avril 2022 - Septembre 2022 :} Postdoc en mathématiques.\\
Collaboration avec Ruadha\'{i} DERVAN.\\
Department of Pure Mathematics and Mathematical Statistics, University of Cambridge, Cambridge, Angleterre.

\bsni \textbf{Septembre 2019 - Décembre 2021 :} Doctorat en mathématiques.\\
\textit{From norms to metrics in non-Archimedean geometry}, co-dirigé par Sébastien BOUCKSOM et Catriona MACLEAN.\\
Institut Fourier, Grenoble.

\bsni \textbf{Septembre 2018 - Juin 2019 :} Master 2 en mathématiques.\\
M2R \textit{Geometry, topology and applications to analysis and probability}.\\
Mémoire de M2 \textit{Multiplier ideals in algebraic geometry}, dirigé par Catriona MACLEAN.\\
Institut Fourier, Grenoble.

\bsni \textbf{Septembre 2017 - Juin 2018 :} Master 1 en mathématiques.\\
Master \textit{Modélisation et analyse numérique}.\\
Mémoire de M1 \textit{Solutions renormalisées de l'équation de transport}, dirigé par Kléber CARRAPATOSO.
Institut Mathématique Alexander Grothendieck, Montpellier.


\bsni \textbf{Septembre 2014 - Juin 2017 :} Licence de mathématiques générales.\\
Licence 2 \textit{CUPGE}, Licence 3 Mathématiques générales.\\
Université d'Avignon et des Pays de Vaucluse, Avignon.

\fi

\par\noindent\rule{\textwidth}{0.4pt}
\textbf{EXPOSES SCIENTIFIQUES.}
%\par\noindent\rule{\textwidth}{0.4pt}


\bsni \textbf{02 Mars 2026} : \textit{Dualité projective et (hyper)déterminants.}\\
\textit{Théminaire}, ENS Lyon.

\bsni \textbf{16 Septembre 2025} : \textit{Autour de l’espace de Goldman–Iwahori.}\\
\textit{Séminaire de Systèmes dynamiques, analyse et géométrie}, Avignon.

\bsni \textbf{15 Mai 2025} : \textit{L’approche variationnelle pour le problème cscK.}\\
\textit{Séminaire GADT}, Dijon.

\bsni \textbf{27 Janvier 2025} : \textit{Arcs and the Hilbert-Mumford criterion for K-stability.}\\
\textbf{Conférence} \textit{Geofloquant}, Bruxelles.

\bsni \textbf{09 Janvier 2025} : \textit{Une version de la conjecture de Yau--Tian--Donaldson via la théorie géométrique des invariants}\\
\textit{Séminaire de géométrie}, Chambéry.

\bsni \textbf{22 Juillet 2024} : \textbf{Mini-cours} sur la géométrie non-archimédienne.\\
\textbf{Conférence} \textit{Non-Archimedean, Tropical and Arakelov Geometry}, Ratisbonne.

\bsni \textbf{01 Juillet 2024} : \textit{Flats in the space of Kähler metrics}.\\
\textbf{Conférence} \textit{BRIDGES meeting on special geometries and stability}, Cargèse.

\bsni \textbf{10 Juin 2024} : \textit{Non-Archimedean methods in special degeneration problems}.\\
\textbf{Conférence} \textit{Freixit (en l'honneur de Gérard Freixas)}, Varberg, Suède.


\bsni \textbf{12-13 Février 2024} : \textbf{Mini-cours }avec Yueqiao Wu sur la stabilité divisorielle, dans le cadre du programme Isaac Newton Institute \textit{New equivariant methods in algebraic and differential geometry}, Cambridge.


\bsni \textbf{25 Janvier 2024} : \textit{Plats dans l'espace des métriques positives}.\\
\textit{Séminaire de géométrie algébrique}, IMJ-PRG, Paris.


\bsni \textbf{11 Janvier 2024} : \textit{Flats in the space of positive metrics.}\\
\textit{KASS Seminar}, Göteborg.


\bsni \textbf{18 Décembre 2023} : \textit{Flats in the space of positive metrics.}\\
\textit{ULB geometry seminar}, Université Libre de Bruxelles.


\bsni \textbf{19 Octobre 2023} : \textit{Transcendental Okounkov bodies.}\\
\textbf{Conférence} \textit{Journées GDR GAGC}, Poitiers.

\bsni \textbf{13 Octobre 2023} : \textit{Corps d'Okounkov transcendents.}\\
\textit{Séminaire de Géométrie, Groupes et Dynamiques}, Institut Camille Jordan, Lyon.

\bsni \textbf{22 Avril 2023} : \textit{The birational geometry of GIT quotients.}\\
\textbf{Conférence} \textit{Around Complex Geometry}, University of Illinois at Chicago.

\bsni \textbf{13 Avril 2023} : \textit{The birational geometry of GIT quotients.}\\
\textbf{Conférence} Oberwolfach \textit{"Komplexe Analysis — Differential and Algebraic methods in Kähler spaces"}, Oberwolfach, exposé court.


\bsni \textbf{04 Avril 2023} : \textit{Stabilité et métriques Kähler–Einstein sur des variétés à fibré anticanonique gros.}\\
\textit{Séminaires systèmes dynamiques et géométrie}, Université d'Angers.

\bsni \textbf{16 Février 2023} : \textit{The birational geometry of GIT quotients.}\\
\textit{KASS Seminar}, Chalmers University of Technology, Göteborg.

\bsni \textbf{20 Janvier 2023} : \textit{Stabilité et métriques Kähler–Einstein sur des variétés à fibré anticanonique gros.}\\
\textit{Séminaire de Géométrie}, Laboratoire de Mathématiques Jean Leray, Nantes.

\bsni \textbf{12 Décembre 2022} : \textit{Stabilité et métriques Kähler–Einstein sur des variétés à fibré anticanonique gros
.}\\
\textit{Séminaire de géométrie complexe}, Institut Elie Cartan de Lorraine.

\bsni \textbf{31 Octobre 2022} : \textit{Ding-stability for manifolds with big anticanonical class.}\\
\textbf{Conférence} \textit{Convex and Complex: Perspectives on Positivity in Geometry. In honor of Bo Berndtsson's 70th birthday}, Cetraro.

\bsni \textbf{25 Mai 2022} : \textit{Geodesics between algebraic models via Berkovich geometry.}\\
\textit{Algebraic Geometry seminar}, DPMMS, Cambridge.

\bsni \textbf{14 Mars 2022} : \textit{Plurisubharmonic geodesics in non-Archimedean geometry and limits of complex manifolds.}\\
\textit{KASS Seminar}, Chalmers University of Technology.

\bsni \textbf{10 Janvier 2022} : \textit{Géodésiques plurisousharmoniques non-archimédiennes et limites de variétés complexes.}\\
\textit{Séminaire d'Analyse Harmonique}, Laboratoire de Mathématiques d'Orsay.

\bsni \textbf{09 Décembre 2021} : \textit{Géodésiques plurisousharmoniques non-archimédiennes et limites de variétés complexes.}\\
\textit{Séminaire de géométrie complexe}, Institut de Mathématiques de Toulouse.

\bsni \textbf{18 Novembre 2021} : \textit{Metric structures on the space of Hermitian norms on a complex vector space.}\\
\textit{Séminaire compréhensible}, Institut Fourier, Grenoble.

\bsni \textbf{08 Novembre 2021} : \textit{Non-Archimedean plurisubharmonic geodesics and limits of complex geodesics.}\\
\textit{Geometric Analysis Seminar}, Purdue.

\bsni \textbf{02 Novembre 2021} : \textit{Legendre transforms, convex bodies, and plurisubharmonic metrics.}\\
\textit{Informal Geometric Analysis Seminar}, University of Maryland.

\bsni \textbf{20 Octobre 2021} : \textit{Plurisubharmonic geodesics in non-Archimedean geometry.}\\
\textit{Oberseminar Algebra und Geometrie}, Frankfurt.

\bsni \textbf{06 Octobre 2021} : \textit{Géodésiques plurisousharmoniques non-archimédiennes et limites de variétés complexes.}\\
\textit{Séminaire de géométrie}, CMLS, Ecole Polytechnique.

\bsni \textbf{07 Septembre 2021} : \textit{Non-Archimedean plurisubharmonic geodesics and complex limits.}\\
\textbf{Conférence} \textit{K-stability and Kähler Geometry}, Cambridge.

\bsni \textbf{09 Février 2021} : \textit{Plurisubharmonic geodesics in non-Archimedean geometry.}\\
\textit{Complex Geometry seminar}, \verb|https://sites.google.com/view/cgseminar|.


\bsni

\par\noindent\rule{\textwidth}{0.4pt}
\textbf{ENSEIGNEMENT ET DIFFUSION.}

\bsni \textbf{Septembre 2024}: Participation à Math@Lyon au collège Jean Jaurès, Villeurbanne.
\\ \verb|https://mmi-lyon.fr/?site_atelier=mathalyon|

\bsni \textbf{2021-2022} - MAT307: Courbes paramétrées et équations différentielles.\\ 
36h TD.\\
\textit{L2 Physique-Mécanique}, DLST, Université Grenoble-Alpes.

\bsni \textbf{2019-2020} - MAT432: Sequences and series of functions.\\
Responsable d'UE, 64h CM+TD.\\
\textit{L2 Math-Info International}, DLST, Université Grenoble-Alpes.

\end{document}
